\section{Testing a Naked Subgraph}
\label{sec:ch08-testing-a-naked-subgraph}
\index{subgraph!testing without a router|(}%

There is no router in this chapter, and that is a choice rather than a schedule.
A subgraph is an ordinary GraphQL service that has agreed to two extra fields,
and both of those fields can be exercised from Postman with nothing else
running. Do that first. Every minute spent learning what a subgraph answers on
its own is a minute you will not spend later staring at a query plan trying to
work out which of three processes is lying.

Two questions, in this order.

\subsection*{What do you publish?}

\index{federation!service field@\code{\_service} field!as a first check}%
\code{\{ \_service \{ sdl \} \}}, to each service, on its own port. That string
is what a composer will read, and reading it yourself takes ten seconds. The
things worth checking are the ones a compiler cannot:

\medskip\noindent\textbf{Is the key there?} \code{@key(fields: "id")} on the
types you meant, and on no others.
Chapter~\ref{ch:how-a-federated-query-actually-runs} showed a schema that built,
started and answered while publishing no key at all, because nothing reachable
from a root field returned the entity.

\medskip\noindent\textbf{Is \code{\_entities} there?} It is generated from the
\code{\_Entity} union, which is generated from the keys. No key, no union, no
field.

\medskip\noindent\textbf{Are the fields on the side you intended?} Mosaic
publishes \code{price}, \code{availableQuantity}, \code{reviews} and
\code{averageRating} on \code{Product}. Catalog publishes the sku, the title,
the description and the category. Both publish \code{id}, and that one is
fine: a key field is shareable by definition, because every subgraph declaring
the entity has to answer with it. Any \emph{other} field on both sides is a
composition error waiting for chapter~\ref{ch:composition}, and a field on
neither is a client error waiting for production.

\subsection*{Can you answer for a key?}

\index{representations!typed by hand}%
Then \code{\_entities}, by hand, exactly as
chapter~\ref{ch:how-a-federated-query-actually-runs} typed it into the sample.
The document is the same one a router would send:

\begin{minted}{graphql}
query ($representations: [_Any!]!) {
  _entities(representations: $representations) {
    ... on Product {
      price { amount currency }
      availableQuantity
      averageRating
    }
  }
}
\end{minted}

\begin{minted}{json}
{ "representations": [ { "__typename": "Product", "id": "UHJvZHVjdDoAAACgAAAAQIAAAAAAAAAB" } ] }
\end{minted}

\begin{minted}[fontsize=\footnotesize]{json}
{"data":{"_entities":[{"price":{"amount":1249.00,"currency":"EUR"},
  "availableQuantity":12,"averageRating":4.25}]}}
\end{minted}

That is Mosaic answering three questions about a product it cannot name. Send
the same key to Catalog and you get the title and the sku and no price. Between
the two responses is the entire supergraph, assembled by hand, and it took two
requests and no new software.

\subsection*{A key that does not decode}

\index{representations!an unresolvable key}%
Send two bad ones on purpose. First, obvious rubbish:

\begin{minted}{json}
{ "representations": [ { "__typename": "Product", "id": "not-a-key" } ] }
\end{minted}

Second, and this is the one that matters, the raw identifier:

\begin{minted}{json}
{ "representations": [ { "__typename": "Product", "id": "a0000000-0000-4000-8000-000000000001" } ] }
\end{minted}

Both answer the same way:

\begin{minted}{json}
{"data":{"_entities":[null]}}
\end{minted}

A null, no \code{errors} key, HTTP 200. Stop on the second one: that string is a
real product identifier, and it is what this field answered before
chapter~\ref{ch:schema-design-that-survives-change} made it a global object
identifier. A representation carrying it is now indistinguishable from garbage,
and the subgraph says so in the mildest possible way.

I think the null is the right answer and I want to be clear about why, because
the first instinct is to throw. \code{\_entities} returns \code{[\_Entity]},
nullable on purpose~\autocite{apollo2026subgraphspec}, so that one bad key in a
batch of twenty-five does not cost the other twenty-four their answers. A
subgraph that throws instead turns a single malformed representation into a
failed fetch, the router turns that into a failed field, and the client gets
nothing for a page of products because one of them was wrong. Fail the entity,
not the batch.

\subsection*{One call for twenty-five, and what happens when it is not}

\index{federation!representation batching!measured}%
Catalog's reference resolver goes through the same DataLoader its node resolver
uses. I asked it for all twenty-five products, first as one \code{\_entities}
call carrying twenty-five representations, then as twenty-five calls carrying
one each, counting the statements that reached PostgreSQL both times:

\begin{center}
\begin{tabular}{lr}
\toprule
How the keys arrived & Statements \\
\midrule
Twenty-five representations, one call & 1 \\
The same twenty-five keys, one per call & 25 \\
\bottomrule
\end{tabular}
\end{center}

\index{DataLoader!behind a reference resolver}%
The data is identical. Nothing in either response says which one happened. This
is chapter~\ref{ch:data-without-the-n-plus-1}'s fix arriving at the door the
router knocks on, and it is why
chapter~\ref{ch:how-a-federated-query-actually-runs} made such a point of
checking that the router sent one call rather than three: the router's batching
and the subgraph's batching are two separate things, and either one of them
being absent costs you the same order of magnitude.

It is worth knowing which half of that you own. The reference resolver itself
does run twenty-five times either way: HotChocolate clones the resolver context
once per representation and starts each one as its own task, which
chapter~\ref{ch:how-a-federated-query-actually-runs} read out of
\code{EntitiesResolver} at tag \code{16.6.0}. So \code{ProductKey.TryDecode}
happens twenty-five times in both rows of that table. What collapses is what
happens next, and it collapses because the resolver hands its key to a
DataLoader rather than to a service. The row that costs twenty-five statements
is the one where the twenty-five resolvers were never in the same request to
begin with, and nothing inside a subgraph can fix that. Whether they arrive
together is the router's decision, which is why
chapter~\ref{ch:how-a-federated-query-actually-runs} spent a section proving
that it batches.

\subsection*{The number that did not move}

\index{measurement!before and after extraction}%
Chapters~\ref{ch:hotchocolate-quickly} to~\ref{ch:schema-design-that-survives-change}
were built on one query and three numbers. The query is gone: no service can
answer it now, because \code{products} belongs to Catalog and \code{reviews}
belongs to Mosaic, and nothing joins them until
chapter~\ref{ch:enter-the-router}. Ask it in two halves instead, twenty-five
products from Catalog and then the same twenty-five keys to Mosaic asking for
each product's reviews and their authors, and this is what Mosaic's timeline
reports:

\begin{minted}[fontsize=\footnotesize,breaklines]{text}
parse - validate 5.019ms compile 2.203ms coerce 3.270ms execute 188.741ms total 204.336ms
    (document cache miss, operation cache miss, 146 resolvers, 2 SQL)
\end{minted}

Ignore the milliseconds. That is the first request a freshly started process
ever answered, both caches missing, so most of the two hundred is the runtime
warming up, exactly as chapter~\ref{ch:the-life-of-a-request} warned. The same
request a moment later reported a total of 4.866 milliseconds with both caches
hitting. It is the two counts at the end of the line that this chapter is
about.

\index{measurement!146 resolvers}%
One hundred and forty-six resolvers, which is exactly what
chapter~\ref{ch:data-without-the-n-plus-1} measured, and it is not a
coincidence. It was one root field, twenty-five review connections and one
hundred and twenty authors. It is now one \code{\_entities} field, twenty-five
review connections and one hundred and twenty authors. The shape of the work did
not change; the name of the field the work hangs from did.

Two statements, where the monolith needed three. The one that left is the
statement that fetched the products, and Catalog runs it instead. Add Catalog's
one to Mosaic's two and you have chapter~\ref{ch:data-without-the-n-plus-1}'s
three, in two processes.

That is the most reassuring measurement in this chapter and I would not read too
much into it. It says the extraction did not cost the database anything, which
is useful and is not the same as saying it cost nothing. What it leaves out is
the round trip between the two services, which does not exist yet because there
is nothing between them.

\subsection*{Verify it in Postman}

\index{Postman!federation collection}%
Every response above is in the companion repository's collection, apart from the
statement counts, which no HTTP assertion can see and which the chapter's
research note records instead. Run the collection rather than reading it. Both
files live in \code{postman/}:

\begin{minted}{text}
mosaic-federation.postman_collection.json
mosaic-federation.local.postman_environment.json
\end{minted}

The environment carries the two URLs. Import both and run the folder:

\begin{minted}{text}
npx newman run postman/mosaic-federation.postman_collection.json \
    -e postman/mosaic-federation.local.postman_environment.json
\end{minted}

\index{Postman!assertions to steal}%
Two assertions in it I would steal for a graph of my own. One sends Catalog
three representations in a deliberately shuffled order and checks that the three
entities come back in that same shuffled order, because the contract is
positional and nothing in the response would tell you if a subgraph had
helpfully sorted them. The other asks Mosaic for \code{products} and asserts
that it fails, which looks like a strange thing to test until you consider what
it catches. A field quietly surviving on the wrong side of a cut is the failure
mode of every extraction, and it is invisible to a suite that only checks that
the right things work.

\index{subgraph!testing without a router|)}%
